\section*{Abstract}
\phantomsection\addcontentsline{toc}{section}{Abstract}
\label{sec:abstract}

% ---- Rework of your draft (your version: 02_abstract_draft.tex).
% ---- Green numbers are from the results pipeline tables.

Solar is one of the fastest-growing sources of electricity, yet a solar string is only as strong as its weakest panel. When a panel in a series string is shaded it limits the current of the whole string, and the bypass diodes used by industry to protect it simply skip the shaded cells and leave their energy on the table. To extract this energy a device coined the `Power Shuffler' has been built: a bidirectional interleaved buck-boost converter that equalises the module voltages by moving current up and down the string. The complexity of modeling the dynamics of the shuffler stems from how physical states are coupled, not just by the coupled inductor itself but also through connection to neighbouring shufflers, so the controller is a multiple-input multiple-output state-space design, developed through a systems-engineering V-model: an analytical model, a high-level Simulink controller architecture, controller component selection, a combined controller-and-hardware simulation in QSPICE, and finally firmware tested on the bench and in the sun. Even in open loop at a fixed 50\,\% duty the shuffler returned power from a partially shaded string. In closed loop, the decoupling PI gain matrix held the module-voltage differences within $\pm0.010\V$ in every settled bench run, below the sensing noise, and simulation shows no overshoot at the deployed gains with 2\,\% settling in 31--41\ms{} at 10\,\% of the simulated gains (18--21\ms{} at 20\,\%). In the solar test, with one of three 80\,W modules shaded to almost zero current, a plain series string would have delivered nothing; the shufflers delivered up to 99\W{} and 10.6\Wh{} over 36 minutes. The controller is stable and extracts power during partial shading events. The next step is a stackable board to test strings of more than three modules.
